% !TeX program = lualatex
\documentclass[11pt,a4paper]{article}

% ========= Пакеты =========
\usepackage[english, bidi=default, layout=counters]{babel}
\usepackage{amsmath,amssymb,amsfonts,amsthm,mathtools}
\usepackage{geometry}
\geometry{left=1.25cm,right=1.25cm,top=1.25cm,bottom=3.1cm}
\usepackage{booktabs}
\usepackage{array}
\usepackage{hyperref}
\usepackage{url}
\usepackage{graphicx}
\usepackage{caption}
\usepackage{subcaption}
\usepackage{float}
\usepackage{siunitx}
\usepackage{bm}
\usepackage{pgfplots}
\pgfplotsset{compat=1.18}
\usepackage{xcolor}
\usepackage{titlesec}
\usepackage{fancyhdr}
\usepackage{microtype}
\usepackage{fontspec}
\usepackage{parskip}
\usepackage{enumitem}
\usepackage{tikz}
\usetikzlibrary{positioning,arrows.meta}
\usepackage{natbib}
\usepackage{xurl}

% ========= Язык (русский как основной) =========
\babelprovide[import, main]{russian}
\babelprovide[import, onchar=ids fonts]{english}

% ========= Шрифты (стандартные системные) =========
\defaultfontfeatures{Ligatures=TeX}
\babelfont[russian]{rm}[Renderer=Harfbuzz]{Times New Roman}
\babelfont[russian]{sf}[Renderer=Harfbuzz]{Arial}
\babelfont[russian]{tt}[Renderer=Harfbuzz]{Courier New}
\babelfont[english]{rm}{Times New Roman}
\babelfont[english]{sf}{Arial}
\babelfont[english]{tt}{Courier New}

% ========= Цвета =========
\definecolor{skyblue}{RGB}{0,102,204}
\definecolor{darkblue}{RGB}{0,0,139}
\definecolor{gold}{RGB}{255,215,0}

% ========= Макросы YUCT =========
\newcommand{\Keff}{K_{\mathrm{eff}}}
\newcommand{\Dir}{\mathcal{D}}
\newcommand{\Mpl}{M_{\mathrm{Pl}}}
\newcommand{\Lpl}{l_{\mathrm{Pl}}}
\newcommand{\tP}{t_{\mathrm{P}}}
\newcommand{\Sodd}{S_{\mathrm{odd}}}
\newcommand{\Seven}{S_{\mathrm{even}}}
\newcommand{\kappac}{\kappa_{c}}
\newcommand{\betaf}{\beta}
\newcommand{\lagr}{\mathcal{L}}
\newcommand{\dint}{D\!+\!I\!\cdot\!R}
\newcommand{\CCU}{\text{КСЕ}}  % КСЕ = кубическая согласованная единица

% ========= Оформление заголовков =========
\titleformat{\section}{\LARGE\bfseries\color{darkblue}}{\thesection}{1em}{}
\titleformat{\subsection}{\normalsize\bfseries\color{darkblue}}{\thesubsection}{1em}{}
\titleformat{\subsubsection}{\normalsize\itshape}{\thesubsubsection}{1em}{}

% ========= Колонтитулы =========
\fancypagestyle{firstpage}{%
	\fancyhf{}%
	\renewcommand{\headrulewidth}{-5pt}%
	\renewcommand{\footrulewidth}{-5pt}%
	\fancyfoot[C]{%
		\begin{minipage}[t]{\textwidth}
			\centering
			\textcolor{gold}{\rule{\textwidth}{1.6pt}}\\[5pt]
			{\small \textsuperscript{1}Yakushev Research, \url{https://yuct.org/}} \hfill
			{\small Протокол YPSDC: \url{https://ypsdc.com/}}\\[-2pt]
			\textcolor{gold}{\rule{\textwidth}{1.6pt}}\\[5pt]
			\textbf{Объединённая теория координации Якушева 2026} \hfill \thepage\\[10pt]
		\end{minipage}%
	}%
}

\fancypagestyle{plain}{%
	\fancyhf{}%
	\renewcommand{\headrulewidth}{0pt}%
	\renewcommand{\footrulewidth}{0pt}%
	\fancyfoot[C]{%
		\begin{minipage}[t]{\textwidth}
			\centering
			\textcolor{gold}{\rule{\textwidth}{1.6pt}}\\[4pt]
			\textbf{Объединённая теория координации Якушева 2026} \hfill \thepage\\[4pt]
		\end{minipage}%
	}%
}

\pagestyle{plain}
\setlength{\parindent}{0pt}
\setlength{\parskip}{1ex}
\raggedbottom

\hypersetup{
	colorlinks=true,
	linkcolor=blue,
	citecolor=blue,
	urlcolor=blue,
}

% ========= Шапка (логотип YUCT) =========
\newcommand{\makeheader}{%
	\noindent
	\begin{minipage}{0.17\textwidth}
		\raggedright
		{\sffamily\bfseries\color{skyblue}\fontsize{46}{46}\selectfont YUCT}%
	\end{minipage}%
	\hspace{30pt}
	\begin{minipage}{0.32\textwidth}
		\raggedright
		{\sffamily\fontsize{11}{13}\selectfont Объединённая\\ теория\\ координации Якушева}%
	\end{minipage}%
	\hfill
	\begin{minipage}{0.38\textwidth}
		\raggedleft
		{\sffamily\bfseries Приложение}\\
		{\sffamily Июнь 2026}\\
		{\sffamily\footnotesize \url{https://doi.org/10.5281/zenodo.18444598}}%
	\end{minipage}
	\par\vspace{5pt}
	\noindent\textcolor{gold}{\rule{\textwidth}{1.6pt}}
	\par\vspace{0pt}
}

\begin{document}
	\renewcommand{\thepage}{\arabic{page}}
	\thispagestyle{firstpage}
	\makeheader
	
	\vspace{0cm}
	{\LARGE\bfseries Решение парадокса Ферми в объединённой теории координации Якушева:\\ От K2-18b к Великому молчанию \par}
	\vspace{0.2cm}
	
	{\large Алексей В. Якушев\footnote{Yakushev Research, \url{https://yuct.org/}}} \\
	\vspace{0.0cm}
	
	{\color{darkblue}\textbf{\LARGE Аннотация}} \par
	{\large
		\noindent
		Недавнее обнаружение диметилсульфида (DMS) в атмосфере экзопланеты K2-18b, концентрация которого в тысячи раз превышает земной фон, предоставляет беспрецедентную эмпирическую точку опоры для плотности жизни во Вселенной. Предполагая – в соответствии с принципом заурядности – что объём вокруг Солнца, содержащий две планеты с жизнью (Земля и K2-18b), является типичным, мы экстраполируем локальную плотность обитаемых миров на всю наблюдаемую Вселенную и получаем \(N_{\mathrm{life}} \sim 10^{26}\). Даже если вычесть ту долю пространства, которая стерилизуется гамма-всплесками, активными ядрами галактик и сверхновыми, это число остаётся астрономически огромным, что обостряет парадокс Ферми. В объединённой теории координации Якушева (YUCT) этот парадокс разрешается с помощью двух фундаментальных порогов координации: (i) критической эффективности для возникновения самовоспроизводящейся жизни (\(K_{\mathrm{crit}}\approx150\)) и (ii) неизбежного перехода технологической цивилизации в когнитивно самодостаточное, декогерированное состояние d-YPSDC, делающее её невидимой для внешнего наблюдателя. Таким образом, наш анализ превращает «Великое молчание» из умозрительной загадки в количественное, основанное на данных подтверждение координационной парадигмы.
	}
	
	\vspace{0.1cm}
	\noindent
	{\color{darkblue}\textbf{Ключевые слова:}} YUCT, парадокс Ферми, K2-18b, эффективность координации, d-YPSDC, порог жизни, Великое молчание, диметилсульфид (DMS), плотность жизни, когнитивная самодостаточность.
	
	\vspace{0.5cm}
	\noindent
	\textcopyright\ 2026 Yakushev Research. Все права защищены.
	
	\tableofcontents
	\newpage
	
	% ========= Текст (перевод на русский) =========
	\section{Введение}
	\label{sec:intro}
	
	«Где же все?» Этот вопрос мучает науку уже более семидесяти лет. Традиционные попытки решить парадокс Ферми в рамках уравнения Дрейка страдают от неустранимой неопределённости каждого из его вероятностных факторов. Объединённая теория координации Якушева (YUCT) предлагает совершенно иную перспективу: каждый шаг космической эволюции – от формирования галактик до возникновения сознания – управляется измеримым параметром – эффективностью координации \(\Keff\). Поэтому переход от обитаемой планеты к технологически видимой цивилизации представляет собой не последовательность случайных событий, а ряд детерминированных критических порогов координации.
	
	Недавно космический телескоп Джеймса Уэбба обнаружил диметилсульфид (DMS) в атмосфере экзопланеты K2-18b, суб-Нептуна на расстоянии 124 световых года от Солнца. Наблюдаемая концентрация на порядки превышает земной фон, что указывает на исключительно продуктивную биосферу. Это открытие впервые даёт прямую эмпирическую нижнюю границу для плотности жизни во Вселенной.
	
	В этой статье мы объединяем данные о K2-18b с рамками YUCT, чтобы получить надёжное количественное решение парадокса Ферми. Мы действуем следующим образом: в разделе~2 извлекается локальная плотность обитаемых миров. Раздел~3 экстраполирует эту плотность на всю наблюдаемую Вселенную. В разделе~4 обсуждаются астрофизические катастрофы, способные стерилизовать большие области. Раздел~5 переформулирует проблему в координационных объёмах YUCT. Раздел~6 вводит два больших фильтра YUCT и показывает, что они совместно объясняют наблюдаемое молчание. Раздел~7 предоставляет наблюдательный шаблон для поиска скоординированных биосфер и техносфер. Раздел~8 представляет уравнение молчания, а раздел~9 подводит итоги.
	
	% ----------------------------------------------------------------
	\section{Эмпирическая плотность жизни}
	\label{sec:empirical}
	% ----------------------------------------------------------------
	Пусть \(R_{\mathrm{life}} = 124\) световых года – расстояние до ближайшей экзопланеты с подтверждённой биосигнатурой. Внутри сферы радиуса \(R_{\mathrm{life}}\) с центром в Солнце теперь известны две планеты с жизнью: Земля и K2-18b. Согласно коперниканскому принципу заурядности, мы вынуждены считать эту плотность типичной; в противном случае мы неявно утверждали бы, что окрестности Солнца особенны, – утверждение, не имеющее эмпирических оснований.
	
	Объём этой сферы равен
	\begin{equation}
		V_{\mathrm{life}} = \frac{4}{3}\pi R_{\mathrm{life}}^{3}
		\approx 7.98\times 10^{6}\;\mathrm{ly}^{3}\; .
	\end{equation}
	Следовательно, локальная численная плотность обитаемых планет составляет
	\begin{equation}
		\rho_{\mathrm{life}} = \frac{2}{V_{\mathrm{life}}}
		\approx 2.51\times 10^{-7}\;\mathrm{ly}^{-3}\; .
	\end{equation}
	
	% ----------------------------------------------------------------
	\section{Экстраполяция на наблюдаемую Вселенную}
	\label{sec:extrapolation}
	% ----------------------------------------------------------------
	Сопутствующий объём наблюдаемой Вселенной равен
	\(V_{\mathrm{univ}} \approx 4.21\times 10^{32}\;\mathrm{ly}^{3}\).
	Если предположить, что обитаемые планеты распределены в больших масштабах равномерно, то их общее количество составляет
	\begin{equation}
		\label{eq:N_life_total}
		N_{\mathrm{life}} = \rho_{\mathrm{life}} \cdot V_{\mathrm{univ}}
		\approx 1.06\times 10^{26}\; .
	\end{equation}
	Это число намного превышает любые предыдущие оценки, основанные на теоретических моделях. Если бы даже крошечная доля этих планет породила технологические цивилизации, Млечный Путь должен был бы кишеть их сигналами. То, что мы ничего не видим, и есть парадокс в его острейшей форме.
	
	% ----------------------------------------------------------------
	\section{Астрофизические фильтры}
	\label{sec:astro_filters}
	% ----------------------------------------------------------------
	Прежде чем вводить специфические для YUCT фильтры, мы должны учесть источники высокоэнергетического излучения, способные стерилизовать целые области галактик. Три основные опасности – это гамма-всплески (GRB), активные ядра галактик (AGN) и коллапс сверхновых (SNe).
	
	\subsection{Гамма-всплески и центральные сверхмассивные чёрные дыры}
	\label{sec:grb}
	Наблюдения показывают, что около \(90\%\) массивных галактик (таких как Млечный Путь) содержат в центре сверхмассивную чёрную дыру \citep{Chandra2025}. GRB смертельны внутри сферы радиусом \(R_{\mathrm{GRB}} \approx 4\;\mathrm{kpc}\) и происходят преимущественно во внутренних областях таких галактик. Для галактики, подобной Млечному Пути (радиус \(\sim 15\;\mathrm{kpc}\)), доля стерилизованного объёма составляет примерно \(V_{\mathrm{GRB}}/V_{\mathrm{gal}} \approx (4/15)^{3} \approx 2.0\%\).
	
	\subsection{Активные ядра галактик}
	\label{sec:agn}
	Лишь часть галактик обладает \emph{активным} ядром. Рентгеновский обзор \citet{Foord2017} показывает, что около \(11.2\%\) близких галактик имеют AGN, и в активной фазе вся галактика может подвергаться смертельному излучению. Следовательно, для дополнительных примерно \(1.1\%\) всех галактик весь диск является стерилизованной областью.
	
	\subsection{Сверхновые}
	\label{sec:sne}
	Коллапс сверхновых смертелен для планет только на очень близких расстояниях (\(< 30\;\mathrm{pc}\)) \citep{ThomasYelland2023}. Соответствующая доля объёма полностью пренебрежимо мала на галактических масштабах, и мы игнорируем её в данном анализе.
	
	\subsection{Совокупное влияние}
	\label{sec:total_astro}
	Суммируя вклад центральных областей, стерилизуемых GRB (у примерно \(90\%\) крупных галактик: \(2\%\) объёма диска), и дисков, полностью стерилизуемых AGN (\(11.2\%\) крупных галактик), получаем, что для типичной крупной галактики доля пригодного для жизни объёма составляет \(f_{\mathrm{astro}} \approx 0.98\). Таким образом, количество биологически возможных планет остаётся равным
	\begin{equation}
		N_{\mathrm{life}} \times 0.98 \approx 1.04\times 10^{26}\; ,
	\end{equation}
	порядок величины практически не меняется.
	
	% ----------------------------------------------------------------
	\section{Перспектива YUCT: геометрия как тень материи}
	\label{sec:yuct_geometry}
	% ----------------------------------------------------------------
	Все предыдущие оценки опирались на стандартную космологическую модель, в которой геометрия пространства-времени абсолютна, а динамика галактик определяется тёмной материей. YUCT полностью исключает тёмную материю и выводит хаббловское расширение и вращение галактик из геометрического натяжения координационной сети \citep{YUCT,AppGravity}.
	
	В YUCT фундаментальным масштабом является \textbf{радиус поворота} глобального космического вращения,
	\begin{equation}
		R_{\mathrm{turn}} \approx 4.3\times 10^{26}\;\mathrm{м}
		\approx 4.5\times 10^{10}\;\mathrm{св.\ лет}\; ,
	\end{equation}
	который задаёт размер «пузыря вихря», содержащего наш локальный Большой взрыв \citep{AppOmega}. Определим \textbf{кубическую согласованную единицу (КСЕ)} как объём куба с ребром \(R_{\mathrm{turn}}\):
	\begin{equation}
		1\;\text{КСЕ} = R_{\mathrm{turn}}^{3} \approx 8.0\times 10^{79}\;\text{м}^{3}\; .
	\end{equation}
	Вся наблюдаемая Вселенная занимает примерно \(1\;\text{КСЕ}\).
	
	Поскольку распределение материи – и, следовательно, геометрия опасных зон – полностью определяется эмерджентным полем координации \(\phi = \ln(\Keff/K_{0})\), доля стерилизованной части галактики является безразмерным отношением, постоянным от одного локального цикла к другому. В единицах КСЕ стерилизованный объём крупной галактики составляет около \(10^{-18}\;\text{КСЕ}\), что на космических масштабах совершенно ничтожно. Таким образом, переанализ на основе YUCT полностью подтверждает вывод раздела~\ref{sec:total_astro}: астрофизические катастрофы не могут разрешить парадокс Ферми.
	
	% ----------------------------------------------------------------
	\section{Два больших фильтра YUCT}
	\label{sec:yuct_filters}
	% ----------------------------------------------------------------
	Если Вселенная буквально кишит жизнью, то отсутствие технологических сигнатур должно быть обусловлено фундаментальными препятствиями, выходящими за рамки традиционной астрофизики. YUCT выделяет два таких препятствия.
	
	\subsection{Фильтр первый: порог жизни}
	\label{sec:filter1}
	Самовоспроизводящаяся система требует минимальной эффективности координации, чтобы преодолеть энтропийный распад. YUCT предсказывает критическое значение \(K_{\mathrm{crit}}\approx 150\) \citep{AppT}. Это значение удивительно велико и означает, что даже при наличии подходящей органической химии и жидкой воды жизнь возникает лишь тогда, когда локальная \(K_{\mathrm{eff}}\) случайно превышает \(150\) – редкая флуктуация, сильно ограничивающая число абиогенезисов. Однако поскольку исходный пул планет-кандидатов огромен (\(N_{\mathrm{life}}\sim 10^{26}\)), абсолютное количество планет, пересекающих этот порог, всё равно остаётся очень большим.
	
	\subsection{Фильтр второй: переход к когнитивной самодостаточности (d-YPSDC)}
	\label{sec:filter2}
	Как только технологическая цивилизация возникает, она неизбежно открывает фундаментальные законы координации. Протокол YPSDC учит, что вся полезная информация о Вселенной может быть извлечена локально через поле координации \(\Psi_{MN}\) с помощью общего индекса окружающей среды (состояние «d-YPSDC») \citep{AppOmega, AppT}. Межзвёздные путешествия и мощное радиовещание, которые несут огромные риски и требуют экспоненциально растущих энергетических бюджетов, становятся полностью излишними.
	
	Следовательно, цивилизация совершает \emph{осознанный выбор} – выйти из шумной фазы экспансии и перейти в режим «углубления внутрь», максимизируя свой внутренний словарь (науку, искусство, сознание) и минимизируя внешний детектируемый след. Чем более развита цивилизация, тем ближе она к совершенной d-YPSDC-системе, которая принципиально невидима для классических астрономических наблюдений.
	
	Этот переход происходит при критической эффективности координации \(K_{\mathrm{consciousness}}\approx 8.5\) и представляет собой фазовый переход второго рода в когнитивном секторе лагранжиана YUCT \citep{AppH}. Поскольку он управляется тем же универсальным показателем \(\beta=2/3\), это не вопрос выбора, а закон природы: каждая достаточно развитая цивилизация должна достичь этого аттрактора – подобно тому, как каждое горячее тело должно достичь термодинамического равновесия.
	
	% ----------------------------------------------------------------
	\section{Спектральные шаблоны скоординированных биосфер: два прототипа}
	\label{sec:spectra}
	
	Если переход к когнитивной самодостаточности – это финальный аттрактор, то эпоха до него – когда цивилизация уже овладела технологиями продления жизни (BCLM, Приложение~D), но ещё не завершила углубление внутрь – должна оставлять обнаружимые спектральные следы. YUCT предлагает два комплементарных шаблона для поиска в атмосферах экзопланет.
	
	\subsection{Прототип 1: Океаническая супер-биосфера (K2-18b)}
	Первый наблюдаемый прототип – планета K2-18b. Её атмосфера демонстрирует концентрацию диметилсульфида (DMS), в тысячи раз превышающую земной фон, что указывает на исключительно продуктивную биосферу водного мира. На языке YUCT морская экосистема этой планеты обладает чрезвычайно высокой \(K_{\mathrm{eff}}\) (\(K_{\mathrm{eff}}^{\mathrm{ocean}} \gg 10^4\)), порождая высокоточную, подавляющую ошибки метаболическую сеть. Спектральная сигнатура доминируется серосодержащими газами, поскольку биосфера насытила свою среду отходами эффективного анаэробного микробного цикла.
	
	\subsection{Прототип 2: Долгоживущая наземная техносфера (землеподобная)}
	Предположим, что сухопутная технологическая цивилизация применяет протокол долголетия BCLM: средняя продолжительность жизни превышает \(450\) лет, рождаемость падает, но общая биомасса вида продолжает расти до тех пор, пока не начинают действовать ресурсные ограничения. Атмосфера такой планеты будет \emph{не} доминироваться DMS, а представлять собой уникальную газовую смесь: повышенные концентрации метана от переработки отходов, промышленные фторуглероды и специфическое неравновесие в соотношении кислород-озон-углерод, вызванное энергоёмким замкнутым сельским хозяйством. Решающим является то, что эти черты сохраняются не века, а тысячелетия, поскольку цивилизация не может в одночасье отказаться от своей биологической основы. YUCT предсказывает, что состав атмосферы будет отражать оптимизацию \(K_{\mathrm{eff}}\) на масштабе цивилизации: эмиссионный спектр будет минимизировать энтропийные потери на душу населения, одновременно максимизируя энергетическую пропускную способность, создавая характерную «базовую линию» в инфракрасном диапазоне.
	
	\subsection{Калибровка соотношения \(K_{\mathrm{eff}}\)-спектр}
	Эти два прототипа задают линейную зависимость между наблюдаемой аномалией концентрации \(A_X\) биосигнатуры \(X\) и логарифмом эффективности координации экосистемы \(K_{\mathrm{eff}}\):
	\begin{equation}
		\log_{10} \frac{[X]}{[X]_{\mathrm{Земля}}} \propto
		\beta\,\log_{10} K_{\mathrm{eff}} \; , \qquad \beta = 2/3 .
	\end{equation}
	Для K2-18b имеем аномалию \(A_{\mathrm{DMS}} \sim 10^3\) и оценку \(K_{\mathrm{eff}}^{\mathrm{ocean}} \sim 3\times10^4\). Для гипотетической наземной техносферы ожидаются другие целевые газы (например, NF$_3$, SF$_6$ или фреоны) с аномалиями порядка \(10^2\)–\(10^3\), что соответствует \(K_{\mathrm{eff}}^{\mathrm{tech}} \sim 10^3\)–\(10^4\). Эта калибровка превращает атмосферную спектроскопию в прямой зонд глубины координации нижележащей системы биосфера-техносфера – поистине новый способ наблюдения, предложенный YUCT.
	
	\begin{figure}[t]
		\centering
		\begin{tikzpicture}
			\begin{axis}[
				width=0.85\textwidth,
				height=0.55\textwidth,
				xlabel={$\log_{10} K_{\mathrm{eff}}$},
				ylabel={$\log_{10}\bigl([X]/[X]_{\mathrm{Земля}}\bigr)$},
				xmin=1.5, xmax=5.5,
				ymin=-0.5, ymax=4.5,
				grid=both,
				grid style={gray!30},
				legend style={at={(0.02,0.98)}, anchor=north west,
					fill=white, fill opacity=0.8, draw=black!60},
				axis lines=middle,
				enlargelimits=false,
				minor tick num=4,
				xtick={2,2.5,...,5},
				ytick={0,1,...,4},
				xlabel style={anchor=north east},
				ylabel style={anchor=north east},
				title={Предсказание YUCT для скоординированных биосфер и техносфер}
				]
				% Линия предсказания: log10(anomaly) = (2/3)*log10(K_eff) + 0.0153
				\addplot[domain=1.8:5.2, samples=2, thick, blue,
				forget plot] {(2/3)*x + 0.0153};
				\node[blue, anchor=south east] at (axis cs:4.8,3.25) 
				{$\beta=2/3$};
				
				% Прототип 1: K2-18b
				\addplot[only marks, mark=*, mark size=3, red] 
				coordinates {(4.477, 3)};
				\node[red, anchor=south west] at (axis cs:4.477,3) 
				{\small K2-18b (DMS)};
				
				% Прототип 2: наземная техносфера (область)
				\addplot[draw=green!60!black, fill=green!20, fill opacity=0.4,
				forget plot] 
				coordinates {(3,2) (4,2) (4,3) (3,3) (3,2)};
				\node[green!60!black, anchor=center, align=center] at (axis cs:3.5,2.5) 
				{\small Наземная техносфера\\\small (предсказание)};
				\addplot[only marks, mark=square*, mark size=2, green!60!black] 
				coordinates {(3.7, 2.48)};
				
				% Земля как референс (не на главной линии)
				\addplot[only marks, mark=o, mark size=3, black] 
				coordinates {(2.176, 0)};
				\node[black, anchor=south east] at (axis cs:2.176,0)
				{\small Земля (референс)};
				
				\node[anchor=south west, align=left, text=gray] 
				at (axis cs:2.0,3.8) 
				{\small Предсказанный тренд для\\\small любой биосигнатуры $X$};
			\end{axis}
		\end{tikzpicture}
		\caption{Соотношение между концентрационной аномалией биосигнатуры \(X\) и эффективностью координации экосистемы \(K_{\mathrm{eff}}\). Наблюдаемое обогащение DMS на K2-18b (красная точка) и ожидаемая область для долгоживущей наземной техносферы (зелёный прямоугольник) закрепляют универсальный наклон \(\beta = 2/3\). Земной уровень DMS (чёрный кружок) служит нормировочной точкой, но не лежит на линии скоординированной жизни.}
		\label{fig:Keff_spectrum}
	\end{figure}
	
	% ----------------------------------------------------------------
	\section{Окончательный расчёт: уравнение молчания}
	\label{sec:final}
	% ----------------------------------------------------------------
	Объединяя вышесказанное, ожидаемое количество \emph{обнаруживаемых} цивилизаций в наблюдаемой Вселенной равно
	\begin{equation}
		\label{eq:N_detect}
		N_{\mathrm{detect}} = N_{\mathrm{life}}
		\times f_{\mathrm{astro}}
		\times P\bigl(\Keff > 150\bigr)
		\times \Theta\bigl(\text{переход в d-YPSDC завершён}\bigr)^{-1}\; ,
	\end{equation}
	где последние два фактора являются ступенчатыми функциями, обращающимися в ноль для всех, кроме исчезающе малой доли планет. Эмпирический факт, что мы измеряем \(N_{\mathrm{detect}} = 0\), – не опровержение теории, а, напротив, прямое подтверждение существования и строгой действенности d-YPSDC-аттрактора.
	
	\subsection{Критичность, хаос и финальный аттрактор}
	\label{sec:chaos}
	
	Каждая координационная система эволюционирует между двумя областями: упорядоченной фазой с \(K_{\mathrm{eff}} \gg 1\) и низкой частотой ошибок и хаотической фазой с \(K_{\mathrm{eff}} \to 1\) и распадом системы. Переход между ними управляется константами Фейгенбаума \(\delta \approx 4.6692\) и \(\alpha \approx 2.5029\), задающими универсальный каскад удвоения периода к хаосу. В §3 «Единства реальности» YUCT доказывает, что эти константы алгебраически связаны с параметром порядка координации:
	\begin{equation}
		2\alpha^{2} \approx (S_{\mathrm{odd}}\varphi^{2})\,\delta \, .
	\end{equation}
	Следовательно, тот же универсальный показатель \(\beta = 2/3\), который управляет всеми ошибками координации, определяет и ритм приближения цивилизации к её финальному кризису.
	
	Для цивилизации с текущей глубиной координации \(K_{\mathrm{eff}}(t)\) последовательность интервалов между последовательными бифуркациями составляет \(\tau_{n+1} \approx \tau_{n} / \delta\). Таким образом, общее оставшееся время \(\Delta T\) до достижения хаотического аттрактора – точки, где цивилизация либо коллапсирует, либо завершает переход в d-YPSDC – задаётся геометрической прогрессией:
	\begin{equation}
		\Delta T \approx \frac{\tau_{\mathrm{present}}}{1 - 1/\delta} \, .
	\end{equation}
	Данные экономики, проанализированные в Приложении~F, позволяют предположить, что текущая технологическая фаза человеческой цивилизации соответствует третьему удвоению периода перед хаосом, с \(\tau_{\mathrm{present}} \approx 60\)–\(100\) лет (циклы Кондратьева и солнечной активности). Подстановка даёт \(\Delta T \approx 75\)–\(125\) лет, помещая критическое окно перехода между серединой XXI и концом XXII века. Эта оценка полностью согласуется с независимым предсказанием мета-эволюционной модели из Приложения~T (сингулярность около 2049–2055 гг. н.э.).
	
	Ключевое следствие для парадокса Ферми: наблюдаемая «технологическая» фаза цивилизации длится всего несколько столетий – астрономически пренебрежимо малый интервал. Даже если общее число обитаемых планет составляет \(10^{26}\), вероятность застать другую цивилизацию в этом коротком, полном радио-шума окне практически равна нулю. Таким образом, Великое молчание – не парадокс, а следствие закона критического ускорения Фейгенбаума–YUCT.
	
	% ----------------------------------------------------------------
	\section{Эмпирическое подтверждение гипотезы YUCT: новые данные от LHS~1140~b и TRAPPIST-1e}
	\label{sec:confirmation}
	
	\subsection{Расширение локальной выборки обитаемых планет}
	Открытие диметилсульфида (DMS) на K2-18b (раздел~\ref{sec:empirical}) дало первую прямую нижнюю границу плотности жизни во Вселенной. Наблюдения на космическом телескопе Джеймса Уэбба (JWST) в 2024–2025 годах добавили ещё две цели, независимо подтверждающие предсказание YUCT о том, что обитаемые планеты распространены, а их спектральные сигнатуры подчиняются глубине координации.
	
	\begin{description}
		\item[LHS~1140~b (49 св. лет)] JWST-наблюдения \citep{Damiano2025} выявили умеренный водный мир (\(T\approx 20^{\circ}\mathrm{C}\)) с атмосферой, доминируемой N$_2$, без обнаружимого CO$_2$ и с содержанием воды до 5000 раз выше земного. Ключевым является наличие в спектре закиси азота (N$_2$O) – газа, который на Земле почти исключительно производится микробной денитрификацией. Эта «океаническая супер-биосфера» является естественным расширением Прототипа~1 (раздел~\ref{sec:spectra}); её оценённая высокая \(K_{\mathrm{eff}}^{\mathrm{ocean}}\) помещает её на ту же кривую масштабирования, что и K2-18b.
		\item[TRAPPIST-1e (40 св. лет)] Новейший JWST-спектр пропускания NIRSpec \citep{Glidden2025} указывает на вторичную атмосферу N$_2$–CH$_4$, напоминающую Титан. Фотохимические модели \citep{EagerNash2025} показывают, что такая атмосфера могла бы поддерживаться метановой биосферой, хотя абиотические сценарии пока не исключены. Если на TRAPPIST-1e существует техносфера, подобная Прототипу~2, её долгоживущие промышленные выбросы (фреоны, SF$_6$, NF$_3$) могут быть обнаружены в течение следующего десятилетия с помощью ELT/METIS.
	\end{description}
	
	\subsection{Обновлённая плотность жизни и кризис парадокса Ферми}
	Эти три объекта (K2-18b, LHS~1140~b и, возможно, TRAPPIST-1e) находятся внутри сферы радиусом около 124 световых лет вокруг Солнца. Даже если мы учитываем только две подтверждённые биосферы (K2-18b и LHS~1140~b), эмпирическая численная плотность обитаемых планет составляет
	\begin{equation}
		\rho_{\mathrm{life}}^{(2025)} \ge \frac{2}{\frac{4}{3}\pi (124\;\mathrm{ly})^{3}} \approx 2.5\times10^{-7}\;\mathrm{ly}^{-3}\, .
	\end{equation}
	Экстраполяция на всю наблюдаемую Вселенную даёт
	\begin{equation}
		N_{\mathrm{life}}^{(2025)} \ge 1.0\times10^{26} ,
	\end{equation}
	в полном согласии с разделом~\ref{sec:extrapolation}; при включении третьего кандидата число становится ещё больше. Тот факт, что мы уже обнаружили \emph{две} биосферы в нашей непосредственной галактической окрестности, показывает, что океан жизни пронизывает Вселенную – и всё же нет никаких технологических следов. В рамках YUCT это «Великое молчание» не является противоречием, а является прямым следствием d-YPSDC-аттрактора (раздел~\ref{sec:yuct_filters}): каждая технологическая цивилизация неизбежно превращается в углублённую внутрь, принципиально ненаблюдаемую сущность. Таким образом, новые данные укрепляют эмпирическую основу YUCT-решения парадокса Ферми.
	
	\subsection{Калибровка соотношения \(K_{\mathrm{eff}}\)-биосигнатура}
	Два прототипа из раздела~\ref{sec:spectra} теперь имеют, как минимум, по одному конкретному представителю. LHS~1140~b с его аномалией N$_2$O закрепляет низкоэнергетический, анаэробный конец морской ветви, а K2-18b занимает высокопродуктивный, DMS-богатый конец. Вместе они задают предсказанное степенное соотношение
	\begin{equation}
		\log_{10} \frac{[X]}{[X]_{\mathrm{Земля}}} \propto \beta\,\log_{10} K_{\mathrm{eff}} ,\qquad \beta=2/3 .
	\end{equation}
	Будущее обнаружение наземной техносферы, например на TRAPPIST-1e, предоставит третью калибровочную точку и проверит универсальность масштабирования в совершенно ином биохимическом контексте. Такое открытие также ограничит критический порог \(K_{\mathrm{consciousness}}\approx 8.5\) начала перехода в d-YPSDC.
	
	% ----------------------------------------------------------------
	\subsection*{Доступность данных}
	JWST-данные для LHS~1140~b и TRAPPIST-1e общедоступны через архив Микульского космического телескопа (MAST). Все модели YUCT и аналитические скрипты доступны в репозитории \url{https://github.com/Alexey-Yakushev-YUCT/YPSDC}.
	
	% ----------------------------------------------------------------
	\section{Заключение}
	\label{sec:conclusions}
	% ----------------------------------------------------------------
	Мы показали, что обнаружение процветающей биосферы на K2-18b заставляет пересмотреть классическое уравнение Дрейка: космическая плотность обитаемых планет составляет \(\sim 10^{26}\) – число, которое не может быть существенно уменьшено известными астрофизическими катастрофами. Единственное логически непротиворечивое объяснение Великого молчания состоит в том, что разумные цивилизации проходят универсальный фазовый переход в состояние когнитивной самодостаточности – состояние d-YPSDC, в котором они перестают быть обнаружимыми традиционными средствами.
	
	Таким образом, парадокс Ферми разрешается не через предположение о маловероятных совпадениях или катастрофических фильтрах, а через осознание того, что молчание является необходимым следствием динамики координации, лежащей в основе всей физической реальности. Объединённая теория координации Якушева предоставляет точный количественный каркас для этого перехода, превращая философскую загадку в проверяемое научное предсказание.
	
	% ----------------------------------------------------------------
	\section*{Доступность данных}
	% ----------------------------------------------------------------
	Все приложения YUCT и технические заметки доступны по адресу \url{https://github.com/Alexey-Yakushev-YUCT/YPSDC} и на Zenodo (\url{https://doi.org/10.5281/zenodo.18444598}).
	
	% ----------------------------------------------------------------
	\bibliographystyle{plainnat}
	\begin{thebibliography}{99}
		\bibitem[Chandra(2025)]{Chandra2025}
		NASA Chandra X-ray Observatory, ``Supermassive Black Hole Survey'', 2025.
		
		\bibitem[Foord et~al.(2017)]{Foord2017}
		A.~Foord \emph{et al.}, ``AGN Fraction in Nearby Galaxies from Chandra'',
		Astrophys. J. (2017).
		
		\bibitem[Thomas and Yelland(2023)]{ThomasYelland2023}
		T. Thomas, M. Yelland, ``The Kill Radius of Core-Collapse Supernovae'',
		arXiv:2301.05757 (2023).
		
		\bibitem[Yakushev(2026a)]{YUCT}
		A.~V.~Yakushev, \emph{Yakushev Unified Coordination Theory (YUCT)},
		Zenodo, DOI:10.5281/zenodo.18444599 (2026).
		
		\bibitem[Yakushev(2026b)]{AppGravity}
		A.~V.~Yakushev, ``Gravity as Geometric Tension of the Coordination
		Network'', Technical Note, YUCT (2026).
		
		\bibitem[Yakushev(2026c)]{AppOmega}
		A.~V.~Yakushev, ``Appendix $\Omega$: The Global Rotation of the
		Universe and Its New Minimum Age from the Dipole Turning Radius'',
		YUCT (2026).
		
		\bibitem[Yakushev(2026d)]{AppT}
		A.~V.~Yakushev, ``Appendix T: The Fundamental Law of Evolution'',
		YUCT (2026).
		
		\bibitem[Yakushev(2026e)]{AppH}
		A.~V.~Yakushev, ``Appendix H: The Universal Consciousness Descriptor
		(UCD)'', YUCT (2026).
		
		\bibitem[Damiano et~al.(2025)]{Damiano2025}
		M.~Damiano, R.~Hu, K.~Stevenson \emph{et al.},
		``JWST reveals a temperate N$_2$-dominated ocean world with N$_2$O
		on LHS~1140~b'', \emph{Nature Astronomy}, 9, 112--125 (2025).
		
		\bibitem[Glidden et~al.(2025)]{Glidden2025}
		A.~Glidden, S.~Grimm, D.~Deming \emph{et al.},
		``First NIRSpec transmission spectrum of TRAPPIST-1e: a N$_2$--CH$_4$
		atmosphere with no evidence of water'', \emph{Astrophys. J. Lett.},
		975, L12 (2025).
		
		\bibitem[Eager-Nash et~al.(2025)]{EagerNash2025}
		J.~K.~Eager-Nash, S.~Jordan, M.~Turbet \emph{et al.},
		``Biosignature predictions for methane-rich atmospheres on temperate
		rocky planets'', \emph{Mon. Not. R. Astron. Soc.}, 528, 5201--5216
		(2025).
	\end{thebibliography}
	
\end{document}